\subsection{Method: NumPy, Tensors, and Vectorization \quad / \quad \textthai{วิธีการ: NumPy, เทนเซอร์ และการประมวลผลแบบเวกเตอร์}}
\begin{paracol}{2}
At the heart of almost every scientific Python project is \texttt{NumPy}. It provides the \texttt{ndarray} object, which allows for fast operations on arrays and matrices. In physics, we often deal with tensors—mathematical objects that describe linear relations between geometric vectors, scalars, and other tensors.

One of \texttt{NumPy}'s most powerful features is \textit{vectorization}\index{Vectorization@การประมวลผลแบบเวกเตอร์ (Vectorization)}. This allows us to perform operations on entire arrays without writing explicit \texttt{for} loops in Python, which are notoriously slow for large datasets.

\switchcolumn

\begin{thai}
หัวใจสำคัญของโครงการ Python เชิงวิทยาศาสตร์เกือบทุกโครงการคือ \texttt{NumPy} ซึ่งจัดเตรียมออบเจ็กต์ \texttt{ndarray} ที่ช่วยให้สามารถดำเนินการกับอาเรย์และเมทริกซ์ได้อย่างรวดเร็ว ในทางฟิสิกส์ เรามักจะจัดการกับเทนเซอร์ (Tensors) ซึ่งเป็นวัตถุทางคณิตศาสตร์ที่อธิบายความสัมพันธ์เชิงเส้นระหว่างเวกเตอร์เรขาคณิต สเกลาร์ และเทนเซอร์อื่นๆ

หนึ่งในฟีเจอร์ที่ทรงพลังที่สุดของ \texttt{NumPy} คือการทำ \textit{vectorization} สิ่งนี้ช่วยให้เราสามารถดำเนินการกับอาเรย์ทั้งหมดได้โดยไม่ต้องเขียนลูป \texttt{for} อย่างชัดเจนใน Python ซึ่งเป็นที่ทราบกันดีว่าช้ามากสำหรับชุดข้อมูลขนาดใหญ่
\end{thai}
\end{paracol}
